\documentclass[10pt,a4paper]{amsart}
\usepackage[margin=19mm]{geometry}
\usepackage{amsmath,amssymb,mathtools}
\usepackage[scr=rsfs,cal=euler]{mathalfa}
\usepackage{xfrac}
\usepackage[unicode,hidelinks,bookmarks=false]{hyperref}
\newcommand{\ie}{i.e.\ }
\newcommand{\cf}{cf.\ }
\newcommand{\resp}{respectively\ }
\newcommand{\Subsection}{\subsection}
\providecommand{\nobreakdash}{\nobreak\hbox{-}}
\setlength{\emergencystretch}{2em}
\multlinegap0pt
\usepackage{mathtools}
\usepackage{float}
\usepackage{amssymb}
\usepackage{tikz}
\usepackage[all]{xy}
\usepackage{xcolor}
\usepackage[shortlabels]{enumitem}
\newtheorem{lemma}{Lemma}
\newcommand{\add}{{\mathrm{add}}}
\title{Derived Picard groups of symmetric representation-finite algebras of type $D$}
\author{Anya Nordskova}
\date{Source: arXiv:2112.09257v2 (CC BY 4.0)}
\begin{document}
\maketitle

\noindent\emph{Source and licence.} Derived Picard groups of symmetric representation-finite algebras of type $D$. Version arXiv:2112.09257v2 (CC BY 4.0), distributed by arXiv under the Creative Commons Attribution 4.0 International licence, http://creativecommons.org/licenses/by/4.0/, as recorded in the arXiv licence metadata for this exact version and shown on the abstract page https://arxiv.org/abs/2112.09257v2. Full licence notice: \texttt{licenses/arxiv-2112.09257-CC-BY-4.0.txt}.\par\smallskip

\noindent\textit{Editorial scope.} Counted unit: the article's lemma thm:twviamut (source0.tex lines 1020--1027). The article's complete original proof of it is delivered with it (source0.tex lines 1029--1050, one \texttt{proof} environment of 22 lines). No mathematics is written, repaired or re-derived: the statement and the proof are the article's own lines, in the article's own notation, and no retained line is substituted or re-flowed.  Retained uncounted as the source-local context the proof needs: lines 968--995.  One declared adaptation: exactly 3 ancillary strings inside the retained spans are omitted (image-inclusion commands and, where the source repeats a label, the repeated label definitions) and no other line differs from the source; the figure environments, with their placement, captions and labels, are kept, and the package records the omitted strings in fidelity receipts bound to the affected bodies.  No retained line contains a citation, so the wrapper carries no bibliography and no bibliography entry.  The retained lines contain the article's own diagram code, which the wrapper typesets with the standard tikz packages; no image file is delivered and the package declares no asset dependency.  Editorial formatting only, in addition: an AMS \texttt{amsart} preamble that restates the article's own declarations and macros used by the retained lines, namely the environments \texttt{lemma} on separate counters, together with the standard packages those lines need.  The retained environments print in this document as Lemma 1 in order of appearance, whereas the article prints the same items with the numbering of its own full document.  The complete native source of the article as distributed in the arXiv e-print for this exact version is bundled unchanged as \texttt{sources/119780/source0.tex}.
\begin{lemma}\label{thm:twonLambda} The spherical twists $t_i$ on $\Lambda$ act on indecomposable projective $\Lambda$-modules in the following way. If $i \geq 4$, 
$$t_i(P_j) =
\begin{cases}
P_j, &\text{ if } j \neq i, i-1\\
P_{j-1},  &\text{ if } j=i\\
(P_i \xrightarrow{\beta_i} P_{i-1} \xrightarrow{soc} P_{i-1})[2],  &\text{ if } j = i-1
\end{cases}
$$ 
If $i=1,2$:

$$
t_i(P_j) =
\begin{cases}
P_j, & \text{ if } j = 3-i\\
P_j[1], & \text{ if } j=i\\
(P_i \xrightarrow{\delta_i\beta_m \dots \beta_{j+1}} P_j)[1], & \text{ if } j \geq 3
\end{cases}
$$ 

Finally, 

$$
t_3(P_j) =
\begin{cases}
(P_3 \xrightarrow{\alpha_1\delta_1\beta_m \dots \beta_{j+1}} P_j)[1], & \text{ if } j \neq 3\\
P_3[1], &\text{ if }  j=3\\
\end{cases}
$$ \end{lemma} 

 \begin{lemma}\label{thm:twviamut}  \begin{enumerate}
     \item  
     $$t_i(\Lambda) \cong (\mu_i^+)^2(\Lambda), t_i^{-1}(\Lambda) = (\mu_{i-1}^-)^2(\Lambda), \, \text{if } i \geq 4.$$
      \item  $$t_3(\Lambda) \cong  \mu_4^- \circ \dots \mu_m^- \circ \mu_2^- \circ \mu_1^-[1](\Lambda)$$
$${t_3^{-1}(\Lambda) \cong  \mu_1^+ \circ  \mu_2^+ \circ \mu_m^+ \dots \mu_4^+ [-1](\Lambda)}$$
     \item $$t_k(\Lambda) \cong  \mu_{3-k}^- \circ \mu_{3}^- \circ \mu_{3-k}^+ \circ \mu_4^- \circ \dots \circ \mu_m^- \circ \mu_{3-k}^-[1](\Lambda)$$ $${t_k^{-1}(\Lambda) \cong  \mu_{3-k}^+ \circ \mu_m^+ \circ  \mu_{3-k}^- \circ \mu_{m-1}^+ \dots \circ \mu_3^+ \circ \mu_{3-k}^+[-1](\Lambda)}, \, \text{if } k = 1,2.$$
    
 \end{enumerate} \end{lemma}

\begin{proof}  The proof is by direct computation. We will only consider the ``positive'' half of the statement and leave the analogous calculations for $t_i^{-1}$'s to the reader.
\begin{enumerate} 
    \item The tilting complex on the left-hand side of $t_i(\Lambda) \cong (\mu_i^+)^2(\Lambda)$ was already computed in Lemma \ref{thm:twonLambda}, so it remains to compute $(\mu_i^+)^2(\Lambda)$. Whenever we have a sequence of mutations applied to $\Lambda$, each subsequent mutation yields a tilting complex and a modified Brauer tree depicting its endomorphism algebra, as discussed in the previous section. There is a natural one-to-one correspondence between the edges of this modified Brauer tree and the indecomposable direct summands of the corresponding tilting complex. Hence, a step-by-step computation of the tilting complex corresponding to a sequence of mutations may be displayed as a sequence of modified Brauer trees whose edges are labeled with complexes. Taking the direct sum of the complexes written on all edges of a modified Brauer tree gives a tilting complex whose endomorphism algebra is depicted by this tree. Every subsequent complex differs from the previous one in only one summand and every subsequent tree differs from the previous one in only one edge. Below we do this for $(\mu_i^+)^2$ and then for the remaining two cases. The curvy red arrows remind the ordering of the relevant edges. Depending of whether we apply $\mu^+$ or $\mu^-$, the arrow points at or away from the edge along which the edge of the next mutation "slides" in the modified Brauer tree. For $(\mu_i^+)^2$ we get the following sequence. 
     \begin{figure}[H]
  
\end{figure}
\noindent Here on the second step we obtained $P_i \xrightarrow{\beta_i} \underline{P_{i-1}}$ as the cone of the minimal left approximation of $P_i$ with respect to $\add(\bigoplus\nolimits_{j\neq i} P_j)$. On the final step we get $P_i \xrightarrow{\beta_i} P_{i-1} \xrightarrow{soc} \underline{P_{i-1}}$ as the cone of the minimal left $\add (\bigoplus\nolimits_{j\neq i} P_j)$-approximation of $P_i \xrightarrow{\beta_i} \underline{P_{i-1}}$. It is easy to see from Lemma \ref{thm:twonLambda} that $t_i(\Lambda) \cong \bigoplus_{j=1}^m t_i(P_j)$ is exactly the complex shown on the right in the picture above.

    \item Now we need to compute the complex $\mu_4^- \circ \dots \mu_m^- \circ \mu_2^- \circ \mu_1^-[1](\Lambda)$. It is convenient to group this sequence this into 2 steps, first applying $\mu_2^- \circ \mu_1^-$ and then the rest of the sequence. 
    
    \begin{figure}[H]
  
\end{figure}

\noindent Same as in the previous case, by Lemma \ref{thm:twonLambda} this is exactly $t_3(\Lambda)$.
    \item The cases $k=1$ and $k=2$ are absolutely analogues, so we assume $k=1$ and compute ${\mu_{2}^- \circ \mu_{3}^- \circ \mu_{2}^+ \circ \mu_4^- \circ \dots \circ \mu_m^- \circ \mu_{2}^-[1](\Lambda)}$. 
     \begin{figure}[H] 
   
  \end{figure}
  \noindent Comparing the result to $t_1(\Lambda)$ obtained in Lemma \ref{thm:twonLambda}, we once again see that this is exactly what was claimed. 
\end{enumerate}
\end{proof} 

\end{document}
